% Blueprint: Character values of finite groups lie in cyclotomic fields 2
% Created by Numina

\begin{document}

% Your blueprint will go here.

\end{document}
% Blueprint: Character values of finite groups lie in cyclotomic fields 2
% Created by Numina

\begin{document}

% Your blueprint will go here.

\end{document}
% Blueprint: Character values of finite groups lie in cyclotomic fields 2
% Created by Numina

\begin{document}

\section{Character values lie in cyclotomic fields}

Let $G$ be a finite group with exponent $n = \exp(G)$. We fix once and for all a
ring embedding $\varphi : \mathbb{Q}(\zeta_n) \hookrightarrow \mathbb{C}$ of the
cyclotomic field, and show that the trace $\operatorname{tr}(\rho(g))$ of every
finite-dimensional complex representation $\rho$ of $G$ at every $g \in G$ lies in
the image $\varphi(\mathbb{Q}(\zeta_n))$.

The argument is elementary and does not use Brauer's induction theorem: since $g^n
= 1$, the operator $\rho(g)$ satisfies $\rho(g)^n = 1$, so each of its eigenvalues
is an $n$-th root of unity. The trace is the sum of these eigenvalues, and the
range of $\varphi$ is a subring of $\mathbb{C}$ containing every $n$-th root of
unity.

\begin{lemma}[Injective embedding of the cyclotomic field]
    \label{lem:cyclotomic-embeds-injective}
    Let $n \neq 0$. There is an injective ring homomorphism
    $\varphi : \mathbb{Q}(\zeta_n) \to \mathbb{C}$.
\end{lemma}
\begin{proof}
    The cyclotomic field $\mathbb{Q}(\zeta_n)$ is a number field, hence algebraic
    over $\mathbb{Q}$, while $\mathbb{C}$ is an algebraically closed
    $\mathbb{Q}$-algebra. By the lifting of embeddings along algebraic extensions
    into an algebraically closed field, there is a $\mathbb{Q}$-algebra
    homomorphism $\mathbb{Q}(\zeta_n) \to \mathbb{C}$; let $\varphi$ be its
    underlying ring homomorphism. A ring homomorphism out of a field is
    injective, so $\varphi$ is injective.
\end{proof}

\begin{lemma}[Embedding carrying a primitive root]
    \label{lem:embedding-primitive-root}
    \uses{lem:cyclotomic-embeds-injective}
    Let $n \neq 0$. There is a ring embedding
    $\varphi : \mathbb{Q}(\zeta_n) \to \mathbb{C}$ and an element
    $\zeta \in \mathbb{C}$ such that $\zeta$ is a primitive $n$-th root of unity
    and $\zeta$ lies in the range of $\varphi$.
\end{lemma}
\begin{proof}
    \uses{lem:cyclotomic-embeds-injective}
    By Lemma~\ref{lem:cyclotomic-embeds-injective} fix an injective ring
    homomorphism $\varphi : \mathbb{Q}(\zeta_n) \to \mathbb{C}$. As a cyclotomic
    extension, $\mathbb{Q}(\zeta_n)$ contains a primitive $n$-th root of unity
    $\xi$. Set $\zeta = \varphi(\xi)$. Because $\varphi$ is injective, transport
    of primitive roots along injective monoid homomorphisms shows $\zeta$ is again
    a primitive $n$-th root of unity, and $\zeta = \varphi(\xi)$ lies in the range
    of $\varphi$.
\end{proof}

\begin{lemma}[Roots of unity lie in a subring containing a primitive root]
    \label{lem:rootsOfUnity-mem}
    Let $n \neq 0$, let $S \subseteq \mathbb{C}$ be a subring, and let
    $\zeta \in S$ be a primitive $n$-th root of unity. Then every $z \in
    \mathbb{C}$ with $z^n = 1$ belongs to $S$.
\end{lemma}
\begin{proof}
    Since $z^n = 1$ with $n \neq 0$, the element $z$ is a unit and lies in the
    group of $n$-th roots of unity of $\mathbb{C}$. As $\zeta$ is a primitive
    $n$-th root of unity, every $n$-th root of unity is a power $\zeta^i$ of
    $\zeta$; in particular $z = \zeta^i$ for some $i$. Because $\zeta \in S$ and
    $S$ is closed under multiplication, $z = \zeta^i \in S$.
\end{proof}

\begin{lemma}[The image of $g$ has finite order]
    \label{lem:rho-pow-exponent}
    Let $\rho$ be a complex representation of $G$ on $V$ and let $g \in G$. Then
    $\rho(g)^{\exp(G)} = 1$ as an endomorphism of $V$.
\end{lemma}
\begin{proof}
    The representation $\rho$ is a monoid homomorphism, so
    $\rho(g)^{\exp(G)} = \rho(g^{\exp(G)})$. By definition of the exponent,
    $g^{\exp(G)} = 1$, and $\rho(1) = 1$. Hence $\rho(g)^{\exp(G)} = 1$.
\end{proof}

\begin{lemma}[Eigenvalues of a finite-order operator are roots of unity]
    \label{lem:eigenvalue-root-of-unity}
    Let $V$ be a finite-dimensional complex vector space, let $T : V \to V$ be a
    linear endomorphism, and let $n \neq 0$ with $T^n = 1$. Then every root $\mu$
    of the characteristic polynomial of $T$ satisfies $\mu^n = 1$.
\end{lemma}
\begin{proof}
    A root of the characteristic polynomial is an eigenvalue of $T$, so $T$ has
    eigenvalue $\mu$. Then $T^n$ has eigenvalue $\mu^n$. But $T^n = 1$ is the
    identity, whose only eigenvalue is $1$ (if $1 \cdot v = \mu^n v$ with
    $v \neq 0$ then $\mu^n = 1$). Therefore $\mu^n = 1$.
\end{proof}

\begin{theorem}[Character values lie in a cyclotomic field]
    \label{thm:brauer-character-in-cyclotomic}
    \uses{lem:embedding-primitive-root, lem:rootsOfUnity-mem, lem:rho-pow-exponent, lem:eigenvalue-root-of-unity}
    Let $G$ be a finite group with exponent $n = \exp(G)$. There is a ring
    embedding $\varphi : \mathbb{Q}(\zeta_n) \to \mathbb{C}$ such that for every
    finite-dimensional complex representation $\rho$ of $G$ on a space $V$ and
    every $g \in G$, the trace $\operatorname{tr}_{\mathbb{C}}(\rho(g))$ lies in
    the range of $\varphi$.
\end{theorem}
\begin{proof}
    \uses{lem:embedding-primitive-root, lem:rootsOfUnity-mem, lem:rho-pow-exponent, lem:eigenvalue-root-of-unity}
    Since $G$ is finite, $n = \exp(G) \neq 0$. By
    Lemma~\ref{lem:embedding-primitive-root} choose an embedding $\varphi$ and a
    primitive $n$-th root of unity $\zeta$ with $\zeta$ in the range of $\varphi$;
    use this $\varphi$. Fix a finite-dimensional representation $\rho$ on $V$ and
    $g \in G$, and write $T = \rho(g)$. By
    Lemma~\ref{lem:rho-pow-exponent}, $T^n = 1$.

    As $\mathbb{C}$ is algebraically closed, the characteristic polynomial of $T$
    splits, so $\operatorname{tr}_{\mathbb{C}}(T)$ equals the sum (with
    multiplicity) of the roots of that polynomial. By
    Lemma~\ref{lem:eigenvalue-root-of-unity}, every such root $r$ satisfies
    $r^n = 1$, and hence by Lemma~\ref{lem:rootsOfUnity-mem} (applied to the
    subring $S = \operatorname{range}(\varphi)$, which contains $\zeta$) each root
    $r$ lies in the range of $\varphi$. The range of $\varphi$ is a subring of
    $\mathbb{C}$, closed under finite sums, so the sum of the roots---namely
    $\operatorname{tr}_{\mathbb{C}}(\rho(g))$---lies in the range of $\varphi$.
\end{proof}

\end{document}
% Blueprint: Character values of finite groups lie in cyclotomic fields 2
% Created by Numina

\begin{document}

\section{Character values lie in cyclotomic fields}

Let $G$ be a finite group with exponent $n = \exp(G)$. We fix once and for all a
ring embedding $\varphi : \mathbb{Q}(\zeta_n) \hookrightarrow \mathbb{C}$ of the
cyclotomic field, and show that the trace $\operatorname{tr}(\rho(g))$ of every
finite-dimensional complex representation $\rho$ of $G$ at every $g \in G$ lies in
the image $\varphi(\mathbb{Q}(\zeta_n))$.

The argument is elementary and does not use Brauer's induction theorem: since $g^n
= 1$, the operator $\rho(g)$ satisfies $\rho(g)^n = 1$, so each of its eigenvalues
is an $n$-th root of unity. The trace is the sum of these eigenvalues, and the
range of $\varphi$ is a subring of $\mathbb{C}$ containing every $n$-th root of
unity.

\begin{lemma}[Embedding carrying a primitive root]
    \label{lem:embedding-primitive-root}
    Let $n \neq 0$. There is a ring embedding
    $\varphi : \mathbb{Q}(\zeta_n) \to \mathbb{C}$ and an element
    $\zeta \in \mathbb{C}$ such that $\zeta$ is a primitive $n$-th root of unity
    and $\zeta$ lies in the range of $\varphi$.
\end{lemma}
\begin{proof}
    The cyclotomic field $\mathbb{Q}(\zeta_n)$ is a number field, hence algebraic
    over $\mathbb{Q}$, while $\mathbb{C}$ is an algebraically closed
    $\mathbb{Q}$-algebra. Therefore there is a $\mathbb{Q}$-algebra homomorphism
    $\mathbb{Q}(\zeta_n) \to \mathbb{C}$ (lifting an embedding along the algebraic
    extension); let $\varphi$ be its underlying ring homomorphism, which is
    injective since it is a homomorphism of fields. As a cyclotomic extension,
    $\mathbb{Q}(\zeta_n)$ contains a primitive $n$-th root of unity $\xi$. Set
    $\zeta = \varphi(\xi)$. Because $\varphi$ is injective, $\zeta$ is again a
    primitive $n$-th root of unity, and $\zeta = \varphi(\xi)$ lies in the range
    of $\varphi$.
\end{proof}

\begin{lemma}[Roots of unity lie in a subring containing a primitive root]
    \label{lem:rootsOfUnity-mem}
    Let $n \neq 0$, let $S \subseteq \mathbb{C}$ be a subring, and let
    $\zeta \in S$ be a primitive $n$-th root of unity. Then every $z \in
    \mathbb{C}$ with $z^n = 1$ belongs to $S$.
\end{lemma}
\begin{proof}
    Since $z^n = 1$ with $n \neq 0$, the element $z$ is a unit and lies in the
    group of $n$-th roots of unity of $\mathbb{C}$. As $\zeta$ is a primitive
    $n$-th root of unity, every $n$-th root of unity is a power $\zeta^i$ of
    $\zeta$; in particular $z = \zeta^i$ for some $i$. Because $\zeta \in S$ and
    $S$ is closed under multiplication, $z = \zeta^i \in S$.
\end{proof}

\begin{lemma}[The image of $g$ has finite order]
    \label{lem:rho-pow-exponent}
    Let $\rho$ be a complex representation of $G$ on $V$ and let $g \in G$. Then
    $\rho(g)^{\exp(G)} = 1$ as an endomorphism of $V$.
\end{lemma}
\begin{proof}
    The representation $\rho$ is a monoid homomorphism, so
    $\rho(g)^{\exp(G)} = \rho(g^{\exp(G)})$. By definition of the exponent,
    $g^{\exp(G)} = 1$, and $\rho(1) = 1$. Hence $\rho(g)^{\exp(G)} = 1$.
\end{proof}

\begin{lemma}[Eigenvalues of a finite-order operator are roots of unity]
    \label{lem:eigenvalue-root-of-unity}
    Let $V$ be a finite-dimensional complex vector space, let $T : V \to V$ be a
    linear endomorphism, and let $n \neq 0$ with $T^n = 1$. Then every root $\mu$
    of the characteristic polynomial of $T$ satisfies $\mu^n = 1$.
\end{lemma}
\begin{proof}
    A root of the characteristic polynomial is an eigenvalue of $T$, so $T$ has
    eigenvalue $\mu$. Then $T^n$ has eigenvalue $\mu^n$. But $T^n = 1$ is the
    identity, whose only eigenvalue is $1$ (if $1 \cdot v = \mu^n v$ with
    $v \neq 0$ then $\mu^n = 1$). Therefore $\mu^n = 1$.
\end{proof}

\begin{theorem}[Character values lie in a cyclotomic field]
    \label{thm:brauer-character-in-cyclotomic}
    \uses{lem:embedding-primitive-root, lem:rootsOfUnity-mem, lem:rho-pow-exponent, lem:eigenvalue-root-of-unity}
    Let $G$ be a finite group with exponent $n = \exp(G)$. There is a ring
    embedding $\varphi : \mathbb{Q}(\zeta_n) \to \mathbb{C}$ such that for every
    finite-dimensional complex representation $\rho$ of $G$ on a space $V$ and
    every $g \in G$, the trace $\operatorname{tr}_{\mathbb{C}}(\rho(g))$ lies in
    the range of $\varphi$.
\end{theorem}
\begin{proof}
    \uses{lem:embedding-primitive-root, lem:rootsOfUnity-mem, lem:rho-pow-exponent, lem:eigenvalue-root-of-unity}
    Since $G$ is finite, $n = \exp(G) \neq 0$. By
    Lemma~\ref{lem:embedding-primitive-root} choose an embedding $\varphi$ and a
    primitive $n$-th root of unity $\zeta$ with $\zeta$ in the range of $\varphi$;
    use this $\varphi$. Fix a finite-dimensional representation $\rho$ on $V$ and
    $g \in G$, and write $T = \rho(g)$. By
    Lemma~\ref{lem:rho-pow-exponent}, $T^n = 1$.

    As $\mathbb{C}$ is algebraically closed, the characteristic polynomial of $T$
    splits, so $\operatorname{tr}_{\mathbb{C}}(T)$ equals the sum (with
    multiplicity) of the roots of that polynomial. By
    Lemma~\ref{lem:eigenvalue-root-of-unity}, every such root $r$ satisfies
    $r^n = 1$, and hence by Lemma~\ref{lem:rootsOfUnity-mem} (applied to the
    subring $S = \operatorname{range}(\varphi)$, which contains $\zeta$) each root
    $r$ lies in the range of $\varphi$. The range of $\varphi$ is a subring of
    $\mathbb{C}$, closed under finite sums, so the sum of the roots---namely
    $\operatorname{tr}_{\mathbb{C}}(\rho(g))$---lies in the range of $\varphi$.
\end{proof}

\end{document}
